In alignment with the Design Science Research Methodology \citep{peffersDesignScienceResearch2007}---which mandates rigorous evaluation to ascertain whether the designed system adequately addresses the specified problem---the Snap and Say application was subjected to empirical testing. The evaluation was structured using the FEDS framework \citep{venableFEDSFrameworkEvaluation2016} to assess the system against the core tension identified during the problem identification phase: the friction-fidelity paradox. The goal was to demonstrate utility, quality, and efficacy \citep{hevnerDesignScienceInformation2004} while validating the proposed agent-driven routing mechanisms.
\label{sec:evaluation}

The system was evaluated through a mixed-methods approach: early simulated benchmarking to quantify baseline nutritional accuracy and autonomous logic efficiency (Objective 1), followed by a human-in-the-loop ``Digital Buffet'' protocol involving older adult participants ($N=19$, Age $\geq 65$) to assess usability and system latency (Objectives 4 and 6). Figure~\ref{fig:user-journey} presents the resulting interface, illustrating the transition from multimodal input capture to the dynamic clarification UI.

\begin{figure}[htbp]
    \centering
    \includegraphics[width=0.65\textwidth]{user_journey.png}
    \caption{Demonstration of the Snap and Say interface mapping multimodal inputs (left) to the dynamic clarification UI (right), reflecting Objectives 4 and 6.}
    \label{fig:user-journey}
\end{figure}

\section{Simulated Benchmarking: Accuracy and Suppression Efficiency}

Prior to human trials, the system's core routing logic (Objective 2) and nutritional estimation accuracy (Objective 1) were evaluated using an automated ``Oracle'' simulation against a stratified $N=500$ subset of the Nutrition5k dataset \citep{thamesNutrition5kAutomaticNutritional2021}, as illustrated in Figure~\ref{fig:evaluation-flow}. The dataset was bifurcated into ``Simple'' items (visually distinct, $\leq 3$ ingredients) and ``Complex'' items (high visual ambiguity or hidden caloric density). The Oracle mechanically supplied ground-truth metadata whenever the agent's Clarification Node generated a question, isolating the system's analytical reasoning capacity from human behavioural variability.

\begin{figure}[htbp]
    \centering
    \includegraphics[width=0.75\textwidth]{evaluation_flow.png}
    \caption{Overview of the evaluation methodology. The automated Oracle protocol (simulated benchmarking) supplies ground-truth metadata to the Clarification Node, isolating the system's analytical reasoning from human variability. The human-in-the-loop Digital Buffet protocol assesses usability with $N=19$ older adults logging standardised meals.}
    \label{fig:evaluation-flow}
\end{figure}

\subsubsection{Estimation Accuracy}
The Agentic Loop was compared against a Single-Shot baseline relying exclusively on a foundational Vision-Language Model (VLM) for zero-shot inference. For Simple items, the Agentic approach yielded a marginal $3.2\%$ reduction in MAE for caloric estimation (74.5\,kcal vs.\ 72.1\,kcal). However, for Complex items, the agentic clarification cycles drove a substantial $30.6\%$ error reduction (185.0\,kcal vs.\ 128.4\,kcal), resulting in an overall MAE improvement of $22.8\%$ across the test bed.

\begin{table}[htbp]
    \centering
    \caption{Nutritional estimation accuracy: Agentic versus Single-Shot baseline.}
    \label{tab:mae-comparison}
    \begin{tabular}{lccc}
        \toprule
        Category & Baseline (Single-Shot) & Snap and Say (Agentic) & Error Reduction \\
        \midrule
        Simple Items  & 74.5\,kcal  & 72.1\,kcal  & 3.2\%  \\
        Complex Items & 185.0\,kcal & 128.4\,kcal & 30.6\% \\
        \textbf{Overall MAE} & \textbf{129.8\,kcal} & \textbf{100.3\,kcal} & \textbf{22.8\%} \\
        \bottomrule
    \end{tabular}
\end{table}

\subsubsection{Suppression Logic Efficiency}
The dynamic complexity scoring architecture successfully bounded unnecessary querying. At the calibrated optimal configuration (Section~\ref{sec:threshold-calibration}), the Clarification Node achieved a $63\%$ True Negative Rate on Simple items (the system autonomously bypassed intervention for unambiguous foods) while correctly triggering intervention on $76\%$ of Complex cases (TPR~$= 0.759$, $95\%$~CI~$[0.655,\, 0.840]$, $N=80$).

\section{Calibration of Routing Thresholds}
\label{sec:threshold-calibration}

Initial benchmarking revealed a performance gap wherein the default configuration yielded a sub-optimal True Negative Rate on Simple items (approximately $54\%$). An empirical threshold tuning experiment was conducted to map the operational space between interaction friction and nutritional fidelity. The parameter sweep evaluated two independent gating mechanisms concurrently:

\begin{enumerate}
    \item \textbf{Deterministic Clinical Threshold ($C_{\text{thresh}}$):} Ranging from 5.0 to 20.0.
    \item \textbf{LLM Confidence Threshold ($\text{Conf}_{\text{thresh}}$):} Ranging from 0.70 to 0.90.
\end{enumerate}

\subsubsection{Stratification Engine Audit}
Prior to executing the parameter sweep, a systematic audit of the Nutrition5k stratification engine was conducted to verify that the Simple/Complex partition was fit for statistical inference. The engine assigns each dish a weighted composite score across four dimensions --- ingredient count ($35\%$), visual distinctiveness ($25\%$), caloric density ($20\%$), and ingredient ambiguity ($20\%$) --- and classifies dishes as Complex when the score exceeds a classification threshold $\theta_s$. Analysis of the full 5,006-dish corpus revealed a heavily left-skewed score distribution (mean~$= 0.194$, median~$= 0.125$, max~$= 0.686$), placing the initial threshold ($\theta_s = 0.50$) at the 97th percentile of the empirical distribution. Consequently, only 185 dishes ($3.7\%$) were classified as Complex, producing a severely imbalanced test set in which a random $N=50$ sample would contain on average fewer than two Complex items --- insufficient to estimate the True Positive Rate with statistical precision. The audit further identified two keyword-level miscalibrations in the heuristic scoring functions. First, the generic term ``mixed'' appeared simultaneously in the set of mixed-dish indicators and in the ambiguity keyword list, causing ingredient names such as ``mixed greens'' --- a visually unambiguous produce item --- to inflate both the visual distinctiveness and ambiguity scores, producing false Complex classifications for dishes with as few as two ingredients. Second, the visual distinctiveness scoring function returned a neutral default of 0.5 for dishes whose ingredient names matched no keywords in either direction; since this value contributed 0.125 to every unrecognised dish's weighted score, it artificially elevated scores for 1,155 dishes ($23\%$) without evidential basis. Three targeted corrections were applied before the sweep was executed: (1) $\theta_s$ was lowered from 0.50 to 0.35, aligning the classification boundary with the 73rd percentile of the score distribution and yielding 1,354 Complex dishes ($27.0\%$); (2) the term ``mixed'' was removed from the mixed-dish indicator set, preventing lexically distinct ingredient names from inflating visual distinctiveness scores; and (3) the visual distinctiveness default was corrected from 0.5 to 0.0, ensuring that the absence of complexity-indicating keywords is interpreted as evidence of simplicity rather than neutral uncertainty.

To obtain reliable estimates across both routing axes, the sweep employed a two-phase design. The first phase ran the full $N=500$ stratified sample across all 30 threshold combinations to characterise the True Negative Rate on Simple items, where the large simple-item pool provided statistically meaningful suppression estimates. The second phase ran a stratified Complex-only sample ($N=80$, seed~$=42$) across six priority combinations ($C_{\text{thresh}} \in \{5.0,\, 8.0\}$, $\text{Conf}_{\text{thresh}} \in \{0.75,\, 0.80,\, 0.85\}$) to obtain precise True Positive Rate and caloric MAE estimates with sufficient Complex-item representation. A combined Pareto analysis across three objectives---maximising TNR, maximising TPR, and minimising caloric MAE---identified $C_{\text{thresh}} = 5.0$, $\text{Conf}_{\text{thresh}} = 0.85$ as the preferred operating point: TPR~$= 0.759$ ($95\%$~CI~$[0.655,\, 0.840]$), TNR~$\approx 0.63$, caloric MAE~$= 92.7\,$kcal. Although this configuration is more aggressive in triggering clarification than a high deterministic threshold alone, it achieved maximum Complex-item information recovery while still suppressing unnecessary intervention on the majority of Simple items. This tuning methodology grounds the system's routing logic in empirical trade-off analysis rather than ad~hoc parameter selection.

\section{User Experience: Digital Buffet Protocol ($N=19$)}
\label{sec:digital-buffet}

Participants were tasked with logging two distinct meal types---one ``Simple'' (roasted chicken) and one ``Complex'' (mixed green salad)---designed to trigger the single-shot and agentic clarification routing paths respectively. Participant responses were collected utilising an adapted survey instrument combining elements of the System Usability Scale (SUS) and NASA Task Load Index \citep{hartDevelopmentNASATLXTask1988}.

Based on the quantitative survey feedback (Q1--Q13 on a 5-point Likert scale):

\begin{enumerate}
    \item \textbf{System Usability and Learnability (Objective 6 Validated).} The overwhelming majority of participants experienced low barriers to entry. Responses indicated strong agreement that the system was easy to learn and that users felt confident utilising the multimodal interface \citep{wildenbosAgingBarriersInfluencing2018}.

    \item \textbf{Perception of Temporal Demand and Complexity.} Scores clustered heavily toward the lower end of the scale (disagreement), indicating a smooth perceived experience despite the complex asynchronous processing.

    \item \textbf{Acceptance of the Clarification Interruption (Objective 4 Validated).} The agentic intervention was generally well-received. Participants largely did not find the system cumbersome or unnecessarily complex.
\end{enumerate}

\section{Qualitative Insights and Edge Cases}

Beyond the numeric metrics, qualitative open-text responses generated critical insights:

\begin{itemize}
    \item \textbf{Latency as Residual Friction.} Several participants explicitly noted speed as an area for improvement. One participant stated, ``De app mag sneller werken'' (The app could work faster). The average 7.1-second round-trip is acceptable for a research prototype \citep{nahStudyTolerableWaiting2004} but requires further optimisation for deployment.

    \item \textbf{Acoustic and System Confidence Challenges.} A subset of users experienced issues with the microphone or speech recognition (e.g., ``Microfoon werkte niet''). When the primary modality failed, the fallback mechanisms became the hardest part of the interaction.

    \item \textbf{Desire for System Competence.} Qualitative feedback requesting features such as ``Let op verschil tussen cal.\ en kcal.'' (Note the difference between cal.\ and kcal.) demonstrates that users quickly escalate their expectations of the system's nutritional expertise once the basic friction of logging is removed.
\end{itemize}

In synthesis, the evaluation confirms the utility of the Snap and Say architecture. The decoupling of unstructured capture from deterministic complexity analysis actively manages the friction-fidelity paradox.
